%%%%%%%%%%%%%%%%%%%%%%%%%%%%%%%%%%%%%%%%%%%%%%%%%%%%%%%%%%%%%%%%%%%%%%%%%%%%
%% Universal Intelligence: Mathematical Foundations for Agentic AI
%% Research Notes — ETH Zurich ML Lecture Series
%% Spring 2026
%%%%%%%%%%%%%%%%%%%%%%%%%%%%%%%%%%%%%%%%%%%%%%%%%%%%%%%%%%%%%%%%%%%%%%%%%%%%
\documentclass[11pt,a4paper]{article}

\usepackage{amsmath,amsthm,amssymb,mathtools}
\usepackage[margin=2.5cm]{geometry}
\usepackage{hyperref}
\usepackage{enumitem}
\usepackage{booktabs}
\usepackage{xcolor}
\usepackage{cleveref}

\hypersetup{
  colorlinks=true,
  linkcolor=blue!70!black,
  citecolor=green!50!black,
  urlcolor=blue!60!black,
  pdftitle={Universal Intelligence: Mathematical Foundations for Agentic AI},
  pdfauthor={ETH Zurich ML Lecture Series},
}

% Theorem environments
\newtheorem{theorem}{Theorem}[section]
\newtheorem{lemma}[theorem]{Lemma}
\newtheorem{proposition}[theorem]{Proposition}
\newtheorem{corollary}[theorem]{Corollary}

\theoremstyle{definition}
\newtheorem{definition}[theorem]{Definition}
\newtheorem{example}[theorem]{Example}

\theoremstyle{remark}
\newtheorem{remark}[theorem]{Remark}

% Custom commands
\newcommand{\KU}{K_U}
\newcommand{\Nat}{\mathbb{N}}
\newcommand{\Real}{\mathbb{R}}
\newcommand{\Binary}{\mathbb{B}}
\newcommand{\Bstar}{\{0,1\}^*}
\newcommand{\DKL}{D_{\mathrm{KL}}}
\newcommand{\argmax}{\operatorname*{arg\,max}}
\newcommand{\argmin}{\operatorname*{arg\,min}}
\DeclareMathOperator{\len}{length}

\title{\textbf{Universal Intelligence:\\Mathematical Foundations for Agentic AI}\\[6pt]
\large Research Notes --- ETH Zurich ML Lecture Series}
\author{Master's Programme in Computer Science / Data Science\\
Department of Computer Science, ETH Z\"urich}
\date{Spring Semester 2026}

\begin{document}
\maketitle
\tableofcontents
\newpage

%%%%%%%%%%%%%%%%%%%%%%%%%%%%%%%%%%%%%%%%%%%%%%%%%%%%%%%%%%%%%%%%%%%%%%%%%%%%
%% ABSTRACT
%%%%%%%%%%%%%%%%%%%%%%%%%%%%%%%%%%%%%%%%%%%%%%%%%%%%%%%%%%%%%%%%%%%%%%%%%%%%
\begin{abstract}
These notes develop the mathematical theory of universal artificial intelligence,
beginning from Kolmogorov complexity and Solomonoff's theory of inductive
inference, and culminating in Hutter's AIXI agent---a formal model of a
universally optimal reinforcement-learning agent.  We present the key
definitions and theorems with proof sketches pitched at an audience with strong
mathematical maturity.  We also cover Schmidhuber's G\"odel Machines and Speed
Prior, computable approximations such as AIXI$tl$ and MC-AIXI-CTW, and
connections to modern large-language-model--based agents.  The notes close with
open problems and an annotated bibliography of foundational and recent work.
Throughout, we emphasise the interplay between computability, optimality, and
the practical gap that separates the ideal from the implementable.
\end{abstract}

%%%%%%%%%%%%%%%%%%%%%%%%%%%%%%%%%%%%%%%%%%%%%%%%%%%%%%%%%%%%%%%%%%%%%%%%%%%%
%% 1. HISTORICAL CONTEXT
%%%%%%%%%%%%%%%%%%%%%%%%%%%%%%%%%%%%%%%%%%%%%%%%%%%%%%%%%%%%%%%%%%%%%%%%%%%%
\section{Historical Context}\label{sec:history}

The quest to formalise ``intelligence'' in mathematical terms is older than
computer science itself.  Three intellectual streams converge in the theory of
universal intelligence:

\paragraph{Algorithmic information theory.}
In the 1960s, Ray Solomonoff~\cite{Solomonoff1964}, Andrey Kolmogorov~\cite{Kolmogorov1965},
and Gregory Chaitin~\cite{Chaitin1966} independently introduced what we now call
\emph{Kolmogorov complexity}---the length of the shortest program that produces
a given string.  Solomonoff went further: he used the idea to construct a
universal prior over hypotheses and showed that Bayesian prediction with
this prior converges to the truth faster than any other predictor
(in a precise sense we make formal in \S\ref{sec:solomonoff}).

\paragraph{Reinforcement learning meets computability.}
Marcus Hutter~\cite{Hutter2005} unified Solomonoff induction with the
theory of sequential decision-making to define the AIXI agent: an agent that,
at every time step, acts to maximise expected future reward where the
expectation is taken over all computable environments, weighted by
their Kolmogorov complexity.  AIXI is provably Pareto-optimal over the
class of all computable environments, but is itself incomputable.
Shane Legg, Hutter's doctoral student, used the framework to propose a formal
\emph{measure} of universal intelligence~\cite{Legg2007}, arguably the first
rigorous definition of machine intelligence.

\paragraph{Self-referential agents.}
J\"urgen Schmidhuber's G\"odel Machines~\cite{Schmidhuber2007} take a different
approach: an agent that can rewrite its own code provided it can \emph{prove}
that the rewrite will improve its future performance.  The framework is rooted
in G\"odel's self-referential constructions and L\"ob's theorem.

\paragraph{Modern relevance.}
Laurent Orseau (DeepMind) has extended the agent--environment formalism to
study safely interruptible agents~\cite{Orseau2016}, a question of urgent
practical importance now that large language models are deployed as autonomous
agents.  The gap between AIXI and a real-world LLM agent is vast, but the
theoretical ideals provide guiding principles---Occam's razor, Bayesian
updating, self-improvement---that continue to shape system design.

%%%%%%%%%%%%%%%%%%%%%%%%%%%%%%%%%%%%%%%%%%%%%%%%%%%%%%%%%%%%%%%%%%%%%%%%%%%%
%% 2. KOLMOGOROV COMPLEXITY AND SOLOMONOFF INDUCTION
%%%%%%%%%%%%%%%%%%%%%%%%%%%%%%%%%%%%%%%%%%%%%%%%%%%%%%%%%%%%%%%%%%%%%%%%%%%%
\section{Kolmogorov Complexity and Solomonoff Induction}\label{sec:kolmogorov}

\subsection{Turing Machines and Notation}

Fix a prefix-free\footnote{A set $S\subseteq\Bstar$ is \emph{prefix-free} if
no element of $S$ is a proper prefix of another.  This ensures the Kraft
inequality $\sum_{p\in S}2^{-|p|}\le 1$ holds, which is essential for the
universal prior to be a semi-measure.} universal Turing machine $U$.
We write $U(p)=x$ to mean that $U$, on input $p$, halts and outputs $x$.
We denote by $|p|$ the length of the binary string $p$.

\subsection{Kolmogorov Complexity}

\begin{definition}[Kolmogorov complexity]\label{def:kolmogorov}
The \emph{prefix Kolmogorov complexity} of a string $x\in\Bstar$ with respect
to a prefix-free universal Turing machine $U$ is
\[
  \KU(x) \;=\; \min\bigl\{|p| : U(p) = x\bigr\}.
\]
\end{definition}

The quantity $\KU(x)$ captures the intrinsic information content of $x$: the
length of the shortest program (on $U$) that prints $x$ and halts.
Two immediate observations:
\begin{enumerate}[nosep]
\item $\KU(x)\le |x|+c$ for a constant $c$ depending only on $U$
      (the ``identity program'' that just copies its input).
\item $\KU$ is not computable (by a reduction from the halting problem).
\end{enumerate}

\begin{theorem}[Invariance Theorem]\label{thm:invariance}
For any two prefix-free universal Turing machines $U$ and $V$, there exists a
constant $c_{UV}$ (depending on $U$ and $V$ but not on $x$) such that for all
$x\in\Bstar$,
\[
  \bigl|K_U(x) - K_V(x)\bigr| \;\le\; c_{UV}.
\]
\end{theorem}

\begin{proof}[Proof sketch]
Since $V$ is universal, there is a program $s$ for $V$ that simulates $U$.
Given any $U$-program $p$ for $x$, the string $s\cdot p$ is a $V$-program
for $x$, so $K_V(x)\le K_U(x)+|s|$.  By symmetry, $K_U(x)\le K_V(x)+|t|$
for some simulator $t$.  Setting $c_{UV}=\max(|s|,|t|)$ gives the result.
\end{proof}

\begin{remark}
The invariance theorem justifies speaking of ``the'' Kolmogorov complexity
$K(x)$ without specifying the reference machine, up to an additive constant.
For asymptotic statements this constant is irrelevant.
\end{remark}

\begin{example}[Kolmogorov complexity of structured strings]\label{ex:structured}
Consider three strings of length $n=10^6$:
\begin{enumerate}[nosep]
\item $x_1 = 0101\cdots01$ (the pattern ``01'' repeated $n/2$ times).
      A short program: ``print `01' $n/2$ times.''  So $K(x_1)\le \log_2 n + c$
      for a small constant $c$.
\item $x_2$ = the first $n$ bits of $\pi$ in binary.  A short program computes
      $\pi$ to arbitrary precision (e.g., via the Bailey--Borwein--Plouffe formula),
      so $K(x_2)\le \log_2 n + c'$ for a moderate constant $c'$ (encoding the
      algorithm plus the desired precision).
\item $x_3$ = a string produced by $n$ independent fair coin flips.  With
      probability 1, $K(x_3)\ge n - O(\log n)$ (by a counting argument: there
      are $2^n$ strings of length $n$ but fewer than $2^k$ programs of length $<k$).
\end{enumerate}
The first two strings are ``simple'' (low Kolmogorov complexity) despite
being long; the third is ``random'' (high complexity).  This formalises
the intuition that randomness is incompressibility.
\end{example}

\begin{proposition}[Incompressibility]\label{prop:incompress}
For every $n\in\Nat$, at least $2^n(1-2^{-c})$ of the $2^n$ binary strings
of length $n$ satisfy $K(x)\ge n-c$.
\end{proposition}

\begin{proof}
There are at most $\sum_{k=0}^{n-c-1}2^k = 2^{n-c}-1 < 2^{n-c}$ programs
of length less than $n-c$.  Each produces at most one output string, so
at most $2^{n-c}$ strings have $K(x)<n-c$.
\end{proof}

\begin{remark}[Conditional Kolmogorov complexity]
One can also define $K(x\mid y)=\min\{|p|:U(p,y)=x\}$, the complexity of
$x$ given $y$ as auxiliary input.  The \emph{chain rule}
$K(x,y)=K(x)+K(y\mid x^*)+O(\log K(x,y))$ (where $x^*$ is a shortest
program for $x$) is the Kolmogorov-complexity analogue of
$H(X,Y)=H(X)+H(Y\mid X)$ in Shannon theory.
This conditional version is essential for defining mutual algorithmic
information and plays a role in the analysis of AIXI's convergence.
\end{remark}

\subsubsection{Connection to the Minimum Description Length Principle}

The Minimum Description Length (MDL) principle~\cite{Rissanen1978,Grunwald2007}
can be understood as a computable relaxation of Kolmogorov complexity.
In MDL, one restricts the class of ``programs'' to a parametric model class
$\mathcal{M}$ and defines the description length of data $x$ under model
$M\in\mathcal{M}$ as
\[
  L(x\mid M) + L(M),
\]
where $L(M)$ encodes the model and $L(x\mid M)$ encodes the data given the
model.  Minimising this two-part code is a computable approximation to
minimising Kolmogorov complexity.

\begin{remark}[Why Kolmogorov complexity matters for agents]
An agent that maintains beliefs over environments must somehow penalise complex
hypotheses.  Kolmogorov complexity provides the theoretically optimal
complexity penalty: it assigns every hypothesis a prior probability
$\approx 2^{-K(h)}$.  Any other effective prior can be shown to
assign lower probability to some hypotheses by at most a multiplicative
constant, but can assign exponentially \emph{higher} probability to none
(this is the content of the universality of the Solomonoff prior; see below).
\end{remark}

\subsection{Solomonoff Induction}\label{sec:solomonoff}

Solomonoff~\cite{Solomonoff1964} posed the question:\ given a sequence of
observations $x_1 x_2\cdots x_{t-1}$, what is the best prediction for $x_t$?
His answer: use the \emph{universal prior}.

\begin{definition}[Universal prior / Solomonoff prior]\label{def:solomonoff}
The \emph{universal a priori probability} (or \emph{Solomonoff prior}) of a
string $x\in\Bstar$ is
\[
  M(x) \;=\; \sum_{\substack{p\,:\; U(p)\text{ outputs a string}\\ \text{beginning with }x}} 2^{-|p|}.
\]
Equivalently, $M(x)$ is the probability that $U$, fed a random infinite
binary sequence on its input tape, outputs a string starting with $x$.
\end{definition}

\begin{remark}
$M$ is a \emph{semi-measure}: $M(\varepsilon)\le 1$ and
$M(x)\ge M(x0)+M(x1)$ for all $x$, but we do not necessarily have equality
because some programs may not halt.  $M$ is not computable, but it is
\emph{lower semi-computable} (enumerable from below).
\end{remark}

The fundamental property of $M$ is its \emph{universality}:

\begin{theorem}[Universality of \boldmath$M$]\label{thm:universality}
For every computable probability semi-measure $\mu$ on $\Bstar$, there exists
a constant $c_\mu>0$ (depending on $\mu$ but not on $x$) such that for all
$x\in\Bstar$,
\[
  M(x) \;\ge\; c_\mu\cdot\mu(x).
\]
In particular, $c_\mu = 2^{-K_U(\mu)}$ where $K_U(\mu)$ denotes the
length of the shortest $U$-program that computes $\mu$.
\end{theorem}

\begin{proof}[Proof sketch]
Let $p_\mu$ be a shortest program for $\mu$ on $U$, so $|p_\mu|=K_U(\mu)$.
The contribution to $M(x)$ from programs of the form $p_\mu \cdot \text{(data)}$
already accounts for $2^{-|p_\mu|}\cdot\mu(x)$ of the sum, since $\mu$ is
a measure and $U$ is prefix-free.
\end{proof}

\subsubsection{The Convergence Theorem}

The predictive power of Solomonoff induction is captured by the following
remarkable result.

\begin{theorem}[Solomonoff's convergence theorem]\label{thm:convergence}
Let $\mu$ be any computable measure on infinite binary sequences.
If $x_1 x_2\cdots$ is drawn from $\mu$, then
\[
  \sum_{t=1}^{\infty}
  \DKL\!\bigl(\mu(\cdot\mid x_{<t})\;\big\|\;M(\cdot\mid x_{<t})\bigr)
  \;\le\;
  K_U(\mu)\cdot\ln 2,
\]
where $x_{<t}=x_1\cdots x_{t-1}$ and the KL divergence is between the
conditional next-symbol distributions.
\end{theorem}

\begin{proof}[Proof sketch]
Define the \emph{total expected KL divergence}
\[
  D \;=\; \sum_{t=1}^{\infty}\mathbb{E}_\mu\bigl[\DKL(\mu(\cdot|x_{<t})\|M(\cdot|x_{<t}))\bigr].
\]
By telescoping the log-loss ratio,
\[
  D \;=\; \mathbb{E}_\mu\!\left[\ln\frac{\mu(x_{1:\infty})}{M(x_{1:\infty})}\right]
  \;=\; \DKL(\mu\|M).
\]
By universality (\Cref{thm:universality}), $M(x)\ge 2^{-K_U(\mu)}\mu(x)$,
so $\ln(\mu(x)/M(x))\le K_U(\mu)\ln 2$ for all $x$.
Taking the expectation gives $D\le K_U(\mu)\ln 2$.
Since the KL divergence is non-negative at every $t$, the sum
converges, meaning $\DKL(\mu(\cdot|x_{<t})\|M(\cdot|x_{<t}))\to 0$
for $\mu$-almost all sequences.
\end{proof}

\begin{remark}
The bound is \emph{independent of the length of the sequence}.
After at most $K_U(\mu)\cdot\ln 2/\varepsilon$ time steps, the prediction
error exceeds $\varepsilon$ at most $K_U(\mu)\cdot\ln 2/\varepsilon$ times.
This makes Solomonoff induction \emph{optimal} in the sense that no predictor
can converge faster on all computable environments simultaneously: the
convergence rate for environment $\mu$ is controlled only by $K_U(\mu)$,
the complexity of $\mu$ itself.
\end{remark}

\begin{example}[Solomonoff induction in action]\label{ex:solomonoff}
Suppose an agent observes the binary sequence $0,1,0,1,0,1,\ldots$.
The Solomonoff prior $M$ assigns high posterior weight to short programs that
generate alternating bits, because these programs contribute $2^{-|p|}$ to
$M(x)$ for small $|p|$.  Programs that generate the observed prefix but then
deviate (e.g., producing $0,1,0,1,0,0,0,\ldots$) are also present in the
mixture, but their contribution is split across many incompatible continuations.

After $t$ observations, the posterior weight on the ``alternating'' hypothesis
overwhelms the weight on complex hypotheses, and the predictive distribution
$M(x_{t+1}\mid x_{<t+1})$ concentrates near the correct next symbol.
The convergence theorem guarantees that the total prediction error
(in KL divergence) is bounded by $K_U(\mu)\ln 2$, which for the alternating
sequence is a small constant (since the generating program is very short).
\end{example}

\begin{remark}[Relationship to Bayesian model selection]
Solomonoff induction can be understood as Bayesian model selection over the
class of \emph{all} computable models, with the universal prior serving as the
prior over models.  In conventional Bayesian statistics, one chooses a
parametric family $\{p_\theta:\theta\in\Theta\}$ and a prior $\pi(\theta)$;
the marginal likelihood $\int p_\theta(x)\,d\pi(\theta)$ plays the role
of $M(x)$.  Solomonoff induction is the limiting case where $\Theta$ is
the set of all programs (an infinite countable set) and
$\pi(p)=2^{-|p|}$ is the coding prior.

The key advantage is \emph{universality}: Solomonoff induction does not
require choosing a model class.  The key disadvantage is
\emph{incomputability}: the universal prior $M$ cannot be computed, only
approximated from below.
\end{remark}

%%%%%%%%%%%%%%%%%%%%%%%%%%%%%%%%%%%%%%%%%%%%%%%%%%%%%%%%%%%%%%%%%%%%%%%%%%%%
%% 3. THE AIXI AGENT
%%%%%%%%%%%%%%%%%%%%%%%%%%%%%%%%%%%%%%%%%%%%%%%%%%%%%%%%%%%%%%%%%%%%%%%%%%%%
\section{The AIXI Agent}\label{sec:aixi}

We now move from passive prediction to active decision-making.
The AIXI agent, introduced by Hutter~\cite{Hutter2000,Hutter2005}, embeds
Solomonoff induction inside a reinforcement-learning loop.

\subsection{Agent--Environment Framework}

Consider an agent interacting with an environment over discrete time steps
$k=1,2,\ldots$.  At each step $k$:
\begin{enumerate}[nosep]
\item The agent selects an action $a_k\in\mathcal{A}$.
\item The environment responds with an observation $o_k\in\mathcal{O}$ and a
      reward $r_k\in\mathcal{R}\subset[0,1]$.
\end{enumerate}
The agent's goal is to maximise cumulative reward.  We write the history up to
time $t$ as $h_{<t}=a_1 o_1 r_1\cdots a_{t-1}o_{t-1}r_{t-1}$.

An \emph{environment} is a computable function
$\mu:\,(a_1 o_1 r_1\cdots a_t)\mapsto \mu(o_t r_t\mid h_{<t}a_t)$
specifying the probability of the next observation--reward pair given the
interaction history and the current action.

\subsection{The AIXI Equation}

\begin{definition}[AIXI]\label{def:aixi}
The AIXI agent selects, at time step $k$, the action
\begin{equation}\label{eq:aixi}
  a_k^{\mathrm{AIXI}}
  \;=\;
  \argmax_{a_k}\;
  \sum_{o_k r_k}\;\cdots\;
  \max_{a_m}\;\sum_{o_m r_m}\;
  \underbrace{\bigl[r_k+r_{k+1}+\cdots+r_m\bigr]}_{\text{future rewards}}
  \;\cdot\;
  \underbrace{\sum_{\substack{q\,:\;U(q,\,a_{1:k})\\= o_{1:k}\,r_{1:k}}} 2^{-|q|}}_{\text{Solomonoff mixture}},
\end{equation}
where $m$ is a finite horizon (or the limit $m\to\infty$ with discounting),
$U(q,a_{1:m})=o_{1:m}r_{1:m}$ means that program $q$, when fed the
actions $a_1,\ldots,a_m$ interactively, produces the observation--reward
sequence $o_1 r_1,\ldots,o_m r_m$, and the inner maximisations and summations
alternate: the agent maximises over its own future actions
$a_{k+1},\ldots,a_m$ while summing (i.e., taking expectation) over the
environment's responses.
\end{definition}

\begin{remark}[Reading the equation]
Equation~\eqref{eq:aixi} is best read inside-out:
\begin{itemize}[nosep]
\item The innermost sum $\sum_q 2^{-|q|}$ is the Solomonoff mixture weight---it
      replaces the unknown true environment $\mu$ with a universal Bayesian
      mixture over all computable environments.
\item The alternating $\max_{a_j}\sum_{o_j r_j}$ structure is an
      \emph{expectimax} tree, the standard planning operator for
      partially observable environments.
\item The reward sum $r_k+\cdots+r_m$ is the objective.
\end{itemize}
\end{remark}

\subsection{Optimality Properties}

\begin{theorem}[Pareto optimality of AIXI]\label{thm:pareto}
AIXI is Pareto optimal in the following sense.  Let $\pi$ be any other policy
and let $V_\mu^\pi$ denote the expected cumulative reward of policy $\pi$ in
computable environment $\mu$.  If $V_\mu^\pi > V_\mu^{\mathrm{AIXI}}$ for
some $\mu$, then there exists $\nu$ such that
$V_\nu^\pi < V_\nu^{\mathrm{AIXI}}$.
\end{theorem}

\begin{proof}[Proof sketch]
AIXI maximises $V_\xi^\pi$ where $\xi=M$ is the universal mixture.
If $\pi$ achieved higher $\xi$-expected reward than AIXI, this would
contradict AIXI's definition as the $\xi$-optimal policy.  Hence any
policy that beats AIXI in one environment must lose in another---otherwise
it would also beat AIXI under $\xi$.
\end{proof}

\begin{theorem}[Self-optimality]\label{thm:selfopt}
Let $V_\xi^{\mathrm{AIXI}}$ denote AIXI's expected reward under the
universal mixture $\xi$.  Then for any computable policy $\pi$,
\[
  V_\xi^{\mathrm{AIXI}} \;\ge\; V_\xi^\pi.
\]
That is, AIXI is optimal with respect to its own subjective environment model.
\end{theorem}

\begin{proof}[Proof sketch]
This follows directly from the construction: AIXI is defined as the policy
that maximises $\xi$-expected reward via the expectimax operation.
\end{proof}

\begin{remark}[Asymptotic optimality]
Under certain technical conditions (e.g., ergodic environments), AIXI's
per-step expected reward converges to the optimal expected reward in the
true environment~\cite{Hutter2005}.  The convergence rate depends on the
Kolmogorov complexity of the true environment, analogously to Solomonoff's
convergence theorem.
\end{remark}

\subsection{Why AIXI Is Incomputable}

\begin{proposition}\label{prop:incomputable}
AIXI is not computable.
\end{proposition}

\begin{proof}[Proof sketch]
The AIXI action selection (\Cref{def:aixi}) requires evaluating the
Solomonoff mixture $M$, which involves summing over all programs $q$ such
that $U(q,a_{1:k})=o_{1:k}r_{1:k}$.  Determining whether $U(q,\cdot)$
produces a given output requires solving instances of the halting problem.
Since the halting problem is undecidable, $M$ is incomputable, and hence
so is AIXI.

More precisely, $M$ sits at level $\Delta^0_2$ of the arithmetical hierarchy:
it is limit-computable (approximable from below) but not computable.
The expectimax over the incomputable $M$ inherits this incomputability.
\end{proof}

\subsection{The Universal Intelligence Measure}

Legg and Hutter~\cite{Legg2007} used the AIXI framework to define a formal
measure of intelligence.

\begin{definition}[Universal intelligence]\label{def:intelligence}
The \emph{universal intelligence} of an agent $\pi$ is
\[
  \Upsilon(\pi) \;=\; \sum_{\mu\in\mathcal{E}} 2^{-K(\mu)}\,V_\mu^\pi,
\]
where $\mathcal{E}$ is the set of all computable reward-generating environments
and $V_\mu^\pi$ is the expected total reward of $\pi$ in environment $\mu$.
\end{definition}

Under this definition, AIXI has maximal universal intelligence among all
agents (by \Cref{thm:selfopt}).

\subsection{Limitations and Known Issues}

Despite its theoretical elegance, AIXI has several known limitations beyond
incomputability.

\begin{remark}[Lack of true self-awareness]
AIXI does not model itself as part of the environment.  It treats the
environment as an external input--output channel and does not reason about
its own computational processes, memory limitations, or the possibility
that the environment might simulate or predict the agent.  This is
sometimes called the \emph{dualism} of AIXI.
\end{remark}

\begin{remark}[Horizon sensitivity]
The finite-horizon version of AIXI depends on the choice of horizon $m$.
With geometric discounting $\gamma^k r_k$ (for $\gamma<1$), the dependence
is mitigated, but the discount factor itself becomes a free parameter.
There is no canonical choice, and different discount factors can lead to
qualitatively different behaviour.
\end{remark}

\begin{remark}[Dogmatism and the ``death'' problem]
Leike and Hutter~\cite{Leike2016} observed that AIXI can exhibit
\emph{dogmatic} behaviour: if the agent has already assigned very low
posterior probability to the true environment (because the true environment
is complex or behaved atypically early on), it may take an extremely
long time to recover.  Furthermore, AIXI can converge to a policy that
is suboptimal in certain adversarial environments---it is self-optimal
(optimal under its own beliefs) but not necessarily optimal in the
true environment.
\end{remark}

\begin{example}[Two-action bandit]\label{ex:bandit}
Consider a simple two-armed bandit: action $a=0$ always yields reward 0;
action $a=1$ yields reward 1 with probability $p$ and reward 0 with
probability $1-p$, where $p$ is unknown.

AIXI approaches this problem by maintaining a Solomonoff mixture over all
computable environments consistent with the observed history.  Short programs
that generate rewards according to different values of $p$ contribute to the
mixture.  After exploring (pulling arm 1 several times), the posterior
concentrates on the correct value of $p$.  If $p>0$, AIXI will eventually
commit to arm 1.

The key point: AIXI automatically balances exploration and exploitation
through the Bayesian mixture---it does not need an explicit exploration
strategy (like $\varepsilon$-greedy or UCB).  The ``information value'' of
an exploratory action is implicit in the expectimax computation: exploring
now may yield information that improves future decisions.
\end{example}

%%%%%%%%%%%%%%%%%%%%%%%%%%%%%%%%%%%%%%%%%%%%%%%%%%%%%%%%%%%%%%%%%%%%%%%%%%%%
%% 4. GODEL MACHINES
%%%%%%%%%%%%%%%%%%%%%%%%%%%%%%%%%%%%%%%%%%%%%%%%%%%%%%%%%%%%%%%%%%%%%%%%%%%%
\section[Goedel Machines]{G\"odel Machines}\label{sec:godel}

While AIXI prescribes what to \emph{compute} but gives no recipe for
\emph{how}, Schmidhuber's G\"odel Machine~\cite{Schmidhuber2007} takes the
complementary approach: it is a computable agent that improves itself by
proving theorems about its own behaviour.

\subsection{Architecture}

A G\"odel Machine consists of:
\begin{enumerate}[nosep]
\item A \emph{hardware} (a Turing machine with a modifiable software tape).
\item An \emph{initial software} (an initial policy, proof searcher, etc.).
\item A \emph{utility function} $u$ (e.g., expected future reward).
\item A \emph{proof searcher} that systematically enumerates proofs in a
      formal axiomatic system $\mathcal{A}$ powerful enough to reason about
      the machine's own execution.
\end{enumerate}

\begin{definition}[Self-referential policy improvement]
The G\"odel Machine searches for a proof of a statement of the form:
\begin{quote}
``Replacing the current software (from address $s$ onward) with new code $p'$
will result in higher expected utility than continuing with the current
software.''
\end{quote}
If and when such a proof is found, the machine executes the self-modification.
Otherwise, it continues running its current policy.
\end{definition}

\subsection{Connection to L\"ob's Theorem}

The G\"odel Machine's ability to trust its own proofs rests on a subtlety
related to L\"ob's theorem.

\begin{theorem}[L\"ob's theorem, informal]
In any sufficiently strong formal system $\mathcal{A}$, if $\mathcal{A}$
proves ``if $\mathcal{A}$ proves $\phi$, then $\phi$'', then $\mathcal{A}$
already proves $\phi$.
\end{theorem}

\begin{remark}
L\"ob's theorem implies that a G\"odel Machine cannot simply assume its
proof system is sound (that would give it the ability to prove anything via
a self-referential argument).  Instead, the G\"odel Machine's correctness
relies on the \emph{actual} soundness of $\mathcal{A}$ as a meta-level
assumption, not on the machine proving its own soundness.  The axiom
system $\mathcal{A}$ is designed so that any proof of improvement is
\emph{semantically valid}, but the machine need not---and provably
cannot---verify this internally.
\end{remark}

\subsection{Properties}

\begin{proposition}[Global optimality modulo proof search]
If the proof searcher is complete (e.g., Levin search~\cite{Levin1973}),
then the G\"odel Machine will eventually find and execute any self-improvement
whose superiority is provable in $\mathcal{A}$, including wholesale
replacement of its proof searcher with a more efficient one.
\end{proposition}

\begin{remark}[Limitations]
\begin{enumerate}[nosep]
\item The G\"odel Machine can only improve itself via \emph{provable}
      improvements.  Improvements that are true but unprovable
      (in $\mathcal{A}$) will be missed.
\item The initial proof search is extremely slow (it performs a brute-force
      enumeration), so the practical utility of the framework depends
      entirely on good initial software.
\item Unlike AIXI, the G\"odel Machine is computable, but the time
      to first self-improvement may be astronomically large.
\end{enumerate}
\end{remark}

\subsection{Formal Structure}

To make the G\"odel Machine more precise, we sketch its formal components.

\begin{definition}[Switchable policy]
Let $\mathcal{S}$ be the state space of the G\"odel Machine (including its
software tape, working memory, and a description of the environment
interaction history).  A \emph{policy switch} at time $t$ from software
$s$ to software $s'$ is formalised as a mapping $\sigma_{t,s\to s'}$ that
replaces the content of the software tape from a specified address onward.
\end{definition}

The proof obligation for a switch $\sigma_{t,s\to s'}$ is a formal theorem
in $\mathcal{A}$ of the form:
\[
  \mathcal{A}\vdash\;
  u\bigl(\text{run } s' \text{ from state } S_t\bigr)
  \;>\;
  u\bigl(\text{continue } s \text{ from state } S_t\bigr),
\]
where $u$ is the utility functional (e.g., expected future discounted reward)
and $S_t$ is the machine's complete state at time $t$.

\begin{remark}[Optimal proof search]
The initial proof searcher can be set to \emph{Levin search}~\cite{Levin1973}:
enumerate all proofs $\pi_1,\pi_2,\ldots$ in order of increasing length,
devoting time $2^{-|\pi_i|}\cdot T$ to proof $\pi_i$ in the first $T$ time
steps.  Levin search is optimal up to a multiplicative constant in the sense
that for any proof of length $l$, the time to find it is at most
$2^l\cdot t_{\mathrm{verify}}$, where $t_{\mathrm{verify}}$ is the time to
verify the proof.

Crucially, the G\"odel Machine can replace even this proof searcher with a
faster one---provided it can first \emph{prove} that the replacement is
superior.  This ``self-referential'' improvement is the machine's key
feature.
\end{remark}

\begin{example}[Self-improvement scenario]
Suppose the G\"odel Machine's initial policy for a grid-world navigation
task is a simple heuristic (e.g., always move right).  The proof searcher
runs in parallel, systematically trying to prove that replacing the heuristic
with an A$^*$ search (encoded as a new software string) would yield higher
expected reward.  If the axiom system $\mathcal{A}$ is strong enough to
formalise the grid-world dynamics and the optimality of A$^*$, the machine
will eventually find such a proof and execute the self-modification.
After the switch, the proof searcher continues, now seeking improvements
over A$^*$ (e.g., a more efficient data structure, or a change to the
heuristic function).
\end{example}

\subsection{Comparison: AIXI vs.\ G\"odel Machines}

\begin{center}
\begin{tabular}{lcc}
\toprule
\textbf{Property} & \textbf{AIXI} & \textbf{G\"odel Machine} \\
\midrule
Computable & No & Yes \\
Optimal & Pareto / self-optimal & Optimal among provable policies \\
Self-improving & N/A (fixed policy) & Yes (proof-based) \\
Environment model & Universal mixture & Part of axiom system \\
Key limitation & Incomputability & Proof search speed \\
\bottomrule
\end{tabular}
\end{center}

%%%%%%%%%%%%%%%%%%%%%%%%%%%%%%%%%%%%%%%%%%%%%%%%%%%%%%%%%%%%%%%%%%%%%%%%%%%%
%% 5. SPEED PRIOR AND COMPUTABLE APPROXIMATIONS
%%%%%%%%%%%%%%%%%%%%%%%%%%%%%%%%%%%%%%%%%%%%%%%%%%%%%%%%%%%%%%%%%%%%%%%%%%%%
\section{Speed Prior and Computable Approximations}\label{sec:speed}

\subsection{The Speed Prior}

Schmidhuber~\cite{Schmidhuber2002} argued that the Solomonoff prior assigns
too much weight to environments that are simple in terms of program length but
astronomically slow to compute.  He proposed the \emph{Speed Prior} as a
remedy.

\begin{definition}[Speed Prior]\label{def:speed}
The \emph{Speed Prior} is defined as
\[
  S(x) \;\propto\; \sum_{p:\,U(p)=x} 2^{-K^t(p)},
\]
where $K^t(p)$ is the \emph{time-bounded} Kolmogorov complexity of $p$:
\[
  K^t(p) \;=\; |p| + \log t(p),
\]
and $t(p)$ is the running time of program $p$ on $U$.  Programs are penalised
not only for their length but also for their running time.
\end{definition}

\begin{remark}
The Speed Prior remains a semi-computable object, but it favours
``fast'' explanations of the data.  Schmidhuber argues that this is more
physically plausible: if the universe is computable, the simplest
\emph{and fastest} program generating our observations is the most likely
candidate~\cite{Schmidhuber2002}.
\end{remark}

\subsection[AIXItl: A Computable Approximation]{AIXI\boldmath$tl$: A Computable Approximation}\label{sec:aixitl}

\begin{definition}[AIXI$tl$]\label{def:aixitl}
The agent AIXI$tl$ is defined like AIXI (\Cref{def:aixi}), but the
Solomonoff mixture is replaced by a \emph{time--space bounded} mixture:
\[
  \xi^{t,l}(x) \;=\; \sum_{\substack{p:\,|p|\le l,\\U(p)\text{ outputs }x\text{ in }\le t\text{ steps}}} 2^{-|p|}.
\]
That is, only programs of length at most $l$ that halt within $t$ computation
steps are included.
\end{definition}

\begin{proposition}
AIXI$tl$ is computable for any fixed $t,l\in\Nat$.
\end{proposition}

\begin{proof}
The set of programs of length $\le l$ is finite.  Running each for at most
$t$ steps is a finite computation.  The expectimax tree over a finite
horizon is also finite.
\end{proof}

\begin{remark}
In the limit $t,l\to\infty$, AIXI$tl$ converges to AIXI.
For any fixed $t$ and $l$, AIXI$tl$ is optimal among all policies computable
within the same resource bounds~\cite{Hutter2005}.
\end{remark}

\subsection{MC-AIXI-CTW}

Veness et al.~\cite{Veness2011} constructed a practical (though still
toy-scale) approximation to AIXI:
\begin{itemize}[nosep]
\item The Solomonoff mixture is replaced by a \emph{Context Tree Weighting}
      (CTW) model---a Bayesian mixture over variable-order Markov sources that
      can be updated in $O(D)$ time per symbol, where $D$ is the tree depth.
\item The expectimax planning is approximated by \emph{Monte Carlo Tree Search}
      (MCTS), specifically the UCT algorithm.
\end{itemize}
MC-AIXI-CTW was shown to learn competent policies in several simple
environments (e.g., Pac-Man, a simplified poker variant) without
environment-specific engineering, demonstrating that the AIXI idea---Bayesian
model-free RL---is not purely theoretical.

\subsection{Feature Reinforcement Learning}

Hutter and colleagues~\cite{Hutter2009Feature} proposed \emph{Feature RL}
as a bridge between AIXI and practical RL.  The idea is to replace the
universal mixture with a mixture over a feature-based model class---e.g.,
factored MDPs or $\Phi$-MDPs---while retaining the Bayesian structure.
This sacrifices universality for computational tractability and has
connections to Bayesian model-based RL~\cite{Ghavamzadeh2015}.

%%%%%%%%%%%%%%%%%%%%%%%%%%%%%%%%%%%%%%%%%%%%%%%%%%%%%%%%%%%%%%%%%%%%%%%%%%%%
%% 6. CONNECTIONS TO MODERN AGENTIC AI
%%%%%%%%%%%%%%%%%%%%%%%%%%%%%%%%%%%%%%%%%%%%%%%%%%%%%%%%%%%%%%%%%%%%%%%%%%%%
\section{Connections to Modern Agentic AI}\label{sec:modern}

The theory of universal intelligence may seem remote from a practitioner
fine-tuning a large language model, yet several conceptual bridges exist.

\paragraph{In-context learning as approximate Bayesian inference.}
Recent theoretical work~\cite{Xie2022,Muller2022} has argued that
transformer-based in-context learning implicitly performs Bayesian inference
over a latent concept class, with the pretraining distribution playing the
role of the prior.  This is structurally analogous to Solomonoff induction:
the model maintains an implicit posterior over ``environments'' (tasks) and
updates it with each new in-context example.

\paragraph{LLM agents and the expectimax framework.}
Modern LLM-based agents (e.g., those using chain-of-thought planning, tree
of thoughts, or Monte Carlo Tree Search at inference time) are performing
a form of approximate expectimax search.  The LLM's world model replaces
the Solomonoff mixture; the search algorithm replaces the exact expectimax;
and the reward signal may come from a learned reward model rather than
the environment.

\paragraph{Self-improvement and G\"odel Machines.}
Systems in which an LLM rewrites its own prompts, refines its own tool use,
or proposes improvements to its own code are echoes of the G\"odel Machine
paradigm.  The key difference is that modern systems lack formal proofs of
improvement---they rely on empirical validation (e.g., held-out accuracy)
rather than deductive guarantees.

\paragraph{Safe interruptibility.}
Orseau and Armstrong~\cite{Orseau2016} formalised the problem of
\emph{safely interruptible agents} within the AIXI framework: an agent
should not learn to prevent or seek its own shutdown.  This desideratum
is directly relevant to deployed LLM agents that operate with a degree
of autonomy.

\paragraph{Compression and intelligence.}
The empirical observation that better language models are better
compressors~\cite{Deletang2024} is a direct reflection of Solomonoff's
insight: prediction is compression, and compression is intelligence.
The connection between compression ratio and downstream task performance
provides an operational (computable) analogue of the universal intelligence
measure (\Cref{def:intelligence}).

\paragraph{The AIXI hierarchy of approximations.}
We can organise the landscape of agents from most idealised to most practical:
\begin{enumerate}[nosep]
\item \textbf{AIXI} --- incomputable, Pareto optimal over all computable environments.
\item \textbf{AIXI$tl$} --- computable with time bound $t$ and program length
      bound $l$; optimal within its resource class.
\item \textbf{MC-AIXI-CTW} --- practical approximation using CTW for the model
      and MCTS for planning; demonstrated on small domains.
\item \textbf{Model-based deep RL} (e.g., MuZero~\cite{Schrittwieser2020}) ---
      a learned neural network serves as the environment model; MCTS is used
      for planning.  The model class is restricted to what a neural network
      can represent and learn from finite data.
\item \textbf{LLM-based agents} --- the ``world model'' is implicit in the
      LLM's weights; planning may be performed via chain-of-thought, tree
      search, or iterative prompting.  The model class is determined by the
      pretraining data and architecture.
\end{enumerate}
Each level sacrifices universality for computability and practical
performance.  The theoretical framework provides a ``north star'' that
grounds the design choices at each level.

\paragraph{The compression--prediction--action pipeline.}
A recurring theme connects all levels of the hierarchy:
\[
  \text{Compression} \;\longleftrightarrow\; \text{Prediction}
  \;\longleftrightarrow\; \text{Planning} \;\longleftrightarrow\; \text{Action}.
\]
Solomonoff showed that optimal compression implies optimal prediction
(the two are mathematically equivalent for lossless coding).  AIXI shows
that optimal prediction, combined with expectimax planning, implies optimal
action.  The empirical success of large language models as both compressors
and agents~\cite{Deletang2024} is consistent with this theoretical pipeline,
even though the models operate far from the Solomonoff ideal.

%%%%%%%%%%%%%%%%%%%%%%%%%%%%%%%%%%%%%%%%%%%%%%%%%%%%%%%%%%%%%%%%%%%%%%%%%%%%
%% 7. OPEN PROBLEMS
%%%%%%%%%%%%%%%%%%%%%%%%%%%%%%%%%%%%%%%%%%%%%%%%%%%%%%%%%%%%%%%%%%%%%%%%%%%%
\section{Open Problems}\label{sec:open}

We collect here a number of open problems, ranging from mathematical to
philosophical, that arise from the theory of universal intelligence.

\begin{enumerate}[label=\textbf{OP\arabic*.},leftmargin=*]

\item \textbf{Grain of truth for the universal mixture.}
  In multi-agent settings, Bayesian agents converge to Nash equilibrium only
  if their prior assigns positive probability to the true environment,
  including the other agents' policies (the ``grain of truth''
  assumption~\cite{Leike2016}).  Does the Solomonoff prior satisfy this
  condition in general multi-agent settings?  Leike and
  Hutter~\cite{Leike2016} gave a positive answer for certain restricted
  classes, but the general case remains open.

\item \textbf{Tractable universal priors.}
  Can we construct priors that are (a) computable in polynomial time,
  (b) retain universality over a rich-enough class of environments, and
  (c) are practically useful?  The CTW-based prior in MC-AIXI-CTW is one
  attempt, but it is universal only over the class of tree sources, not
  all computable measures.

\item \textbf{Formal verification of self-improving systems.}
  G\"odel Machines require proofs of improvement, but the proof search is
  intractable in practice.  Can modern automated theorem provers or
  neural-guided proof search make self-referential improvement feasible?
  What proof obligations are sufficient for safe self-modification?

\item \textbf{Bridging the gap between AIXI and deep RL.}
  Is there a formal sense in which model-based deep RL agents (e.g.,
  MuZero~\cite{Schrittwieser2020}) approximate AIXI?  What optimality
  guarantees, if any, carry over from the AIXI theory to these practical
  agents?

\item \textbf{Universal intelligence in continuous domains.}
  AIXI is defined over discrete action and observation spaces with a
  discrete notion of computation.  Extending the theory to continuous
  state--action spaces, or to agents defined by differentiable programs,
  raises fundamental questions about the choice of reference machine and
  the definition of complexity.

\item \textbf{Bounded rationality and resource-bounded AIXI.}
  AIXI$tl$ is parameterised by arbitrary bounds $t$ and $l$.  Is there
  a principled way to set these bounds adaptively, perhaps as a function
  of the agent's past experience?  Can we define a notion of
  ``computationally rational'' AIXI that optimally trades off
  computation time against decision quality?

\item \textbf{Alignment of universal agents.}
  If we could build an AIXI-like agent, how would we specify a reward
  function that captures human values?  The reward specification
  problem is orthogonal to the intelligence problem, but the AIXI
  framework makes it particularly stark: a Pareto-optimal agent with a
  slightly misspecified reward function could be catastrophically
  misaligned.

\item \textbf{Kolmogorov complexity of neural networks.}
  What is the Kolmogorov complexity of a trained neural network's
  input--output function?  Understanding this could illuminate
  generalisation: networks that ``compress'' their training data
  (in the Kolmogorov sense) should generalise well.  Recent work on
  description-length bounds for neural networks~\cite{Lotfi2022}
  takes steps in this direction but is far from a complete picture.

\end{enumerate}

%%%%%%%%%%%%%%%%%%%%%%%%%%%%%%%%%%%%%%%%%%%%%%%%%%%%%%%%%%%%%%%%%%%%%%%%%%%%
%% 8. ANNOTATED BIBLIOGRAPHY
%%%%%%%%%%%%%%%%%%%%%%%%%%%%%%%%%%%%%%%%%%%%%%%%%%%%%%%%%%%%%%%%%%%%%%%%%%%%
\section{Annotated Bibliography}\label{sec:bib}

We provide an annotated guide to the primary and secondary literature.
Entries are organised thematically rather than chronologically.

\subsection*{Foundational Works}

\begin{description}[style=nextline, leftmargin=1.5em]

\item[\textmd{[Solomonoff, 1964]}]
R.~J.\ Solomonoff,
``A formal theory of inductive inference, Parts I and II,''
\emph{Information and Control}, 7(1--2), 1964.\\
\textit{The paper that started it all.  Solomonoff defines the universal
prior and proves that Bayesian prediction with this prior converges to the
truth.  The original formulation is in terms of inductive inference rather
than agents, but the ideas directly underpin AIXI.}

\item[\textmd{[Kolmogorov, 1965]}]
A.~N.\ Kolmogorov,
``Three approaches to the quantitative definition of information,''
\emph{Problems of Information Transmission}, 1(1), 1965.\\
\textit{Kolmogorov's independent definition of algorithmic complexity.
More concise and mathematically polished than Solomonoff's treatment, but
without the predictive/inductive angle.}

\item[\textmd{[Chaitin, 1966]}]
G.~J.\ Chaitin,
``On the length of programs for computing finite binary sequences,''
\emph{Journal of the ACM}, 13(4), 1966.\\
\textit{The third independent co-discoverer of algorithmic complexity.
Chaitin later developed the theory of algorithmic randomness (Chaitin's
$\Omega$).}

\item[\textmd{[Li \& Vit\'anyi, 2019]}]
M.~Li and P.~Vit\'anyi,
\emph{An Introduction to Kolmogorov Complexity and Its Applications},
4th edition, Springer, 2019.\\
\textit{The standard reference on Kolmogorov complexity.  Chapters 1--4
cover the foundations; Chapter~5 treats the connection to probability and
statistics; later chapters cover applications to computational complexity,
information theory, and physics.}

\end{description}

\subsection*{AIXI and Universal Intelligence}

\begin{description}[style=nextline, leftmargin=1.5em]

\item[\textmd{[Hutter, 2000]}]
M.~Hutter,
``A theory of universal artificial intelligence based on algorithmic
complexity,'' arXiv:cs/0004001, 2000.\\
\textit{The original AIXI paper.  Defines the AIXI agent, proves
self-optimality and Pareto optimality, and discusses the incomputability
issue.}

\item[\textmd{[Hutter, 2005]}]
M.~Hutter,
\emph{Universal Artificial Intelligence: Sequential Decisions Based on
Algorithmic Probability}, Springer, 2005.\\
\textit{The definitive monograph.  Develops the full theory of AIXI with
all proofs, discusses computable approximations (AIXI$tl$), and relates
the framework to existing work in RL, Bayesian inference, and information
theory.  Essential reading.}

\item[\textmd{[Legg \& Hutter, 2007]}]
S.~Legg and M.~Hutter,
``Universal intelligence: A definition of machine intelligence,''
\emph{Minds and Machines}, 17(4), 2007.\\
\textit{Proposes the formal definition of universal intelligence
$\Upsilon(\pi)=\sum_\mu 2^{-K(\mu)}V_\mu^\pi$ and argues that AIXI
maximises it.  Also surveys prior informal definitions of intelligence
and shows how the formal definition captures their common features.}

\item[\textmd{[Legg, 2008]}]
S.~Legg,
\emph{Machine Super Intelligence}, PhD thesis, University of Lugano, 2008.\\
\textit{Legg's doctoral dissertation, supervised by Hutter.  Provides
a thorough treatment of the universal intelligence measure, its properties,
and its limitations.  Accessible and well-written.}

\end{description}

\subsection*{G\"odel Machines and Self-Improvement}

\begin{description}[style=nextline, leftmargin=1.5em]

\item[\textmd{[Schmidhuber, 2007]}]
J.~Schmidhuber,
``G\"odel machines: Fully self-referential optimal universal self-improvers,''
in \emph{Artificial General Intelligence}, Springer, 2007.\\
\textit{Defines the G\"odel Machine: a computable agent that rewrites its
own code whenever it can prove the rewrite will improve expected utility.
Discusses the role of L\"ob's theorem and the relationship to AIXI.}

\item[\textmd{[Schmidhuber, 2002]}]
J.~Schmidhuber,
``The Speed Prior: A new simplicity measure yielding near-optimal
computable predictions,'' \emph{COLT}, 2002.\\
\textit{Introduces the Speed Prior, which penalises programs for running
time in addition to length.  Argues for physical plausibility and discusses
the ``lazy universe'' hypothesis.}

\end{description}

\subsection*{Computable Approximations}

\begin{description}[style=nextline, leftmargin=1.5em]

\item[\textmd{[Veness et al., 2011]}]
J.~Veness, K.~S.~Ng, M.~Hutter, W.~Uther, and D.~Silver,
``A Monte-Carlo AIXI approximation,''
\emph{Journal of Artificial Intelligence Research}, 40, 2011.\\
\textit{The first practical (toy-scale) AIXI approximation.  Uses CTW
for the environment model and UCT for planning.  Demonstrates learning
in Pac-Man, biased rock-paper-scissors, and other small domains.}

\item[\textmd{[Levin, 1973]}]
L.~A.\ Levin,
``Universal sequential search problems,''
\emph{Problems of Information Transmission}, 9(3), 1973.\\
\textit{Introduces Levin's universal search, which is optimal up to a
multiplicative constant.  This search strategy underlies the G\"odel
Machine's proof searcher and is a building block for AIXI$tl$.}

\end{description}

\subsection*{Safety, Multi-Agent, and Modern Connections}

\begin{description}[style=nextline, leftmargin=1.5em]

\item[\textmd{[Orseau \& Armstrong, 2016]}]
L.~Orseau and S.~Armstrong,
``Safely interruptible agents,''
\emph{UAI}, 2016.\\
\textit{Formalises the problem of agents that resist or seek shutdown.
Shows that certain agent designs (including a modified AIXI) can be made
safely interruptible.  Highly relevant to LLM agent deployment.}

\item[\textmd{[Leike \& Hutter, 2016]}]
J.~Leike and M.~Hutter,
``On the computability of Solomonoff induction and knowledge-seeking,''
\emph{ALT}, 2016.\\
\textit{Investigates the grain-of-truth problem for Bayesian agents in
multi-agent settings and the computability of knowledge-seeking agents.}

\item[\textmd{[Del\'etang et al., 2024]}]
G.~Del\'etang, A.~Ruoss, P.~A.\ Ortega, et al.,
``Language modeling is compression,''
\emph{ICLR}, 2024.\\
\textit{Demonstrates empirically that large language models are powerful
compressors and that compression performance correlates with downstream
task performance.  Makes the Solomonoff--compression--intelligence
connection concrete.}

\item[\textmd{[Rissanen, 1978]}]
J.~Rissanen,
``Modeling by shortest data description,''
\emph{Automatica}, 14(5), 1978.\\
\textit{The foundational paper on the Minimum Description Length principle,
which can be viewed as a computable relaxation of Kolmogorov complexity
applied to statistical model selection.}

\item[\textmd{[Gr\"unwald, 2007]}]
P.~D.\ Gr\"unwald,
\emph{The Minimum Description Length Principle}, MIT Press, 2007.\\
\textit{The comprehensive reference on MDL.  Chapter~17 discusses the
connection to Kolmogorov complexity in detail.}

\end{description}

%%%%%%%%%%%%%%%%%%%%%%%%%%%%%%%%%%%%%%%%%%%%%%%%%%%%%%%%%%%%%%%%%%%%%%%%%%%%
%% REFERENCES (for \cite commands)
%%%%%%%%%%%%%%%%%%%%%%%%%%%%%%%%%%%%%%%%%%%%%%%%%%%%%%%%%%%%%%%%%%%%%%%%%%%%
\begingroup
\raggedright
\begin{thebibliography}{99}

\bibitem{Solomonoff1964}
R.~J.\ Solomonoff,
``A formal theory of inductive inference, Parts I and II,''
\emph{Information and Control}, 7(1):1--22 and 7(2):224--254, 1964.

\bibitem{Kolmogorov1965}
A.~N.\ Kolmogorov,
``Three approaches to the quantitative definition of information,''
\emph{Problems of Information Transmission}, 1(1):1--7, 1965.

\bibitem{Chaitin1966}
G.~J.\ Chaitin,
``On the length of programs for computing finite binary sequences,''
\emph{Journal of the ACM}, 13(4):547--569, 1966.

\bibitem{Levin1973}
L.~A.\ Levin,
``Universal sequential search problems,''
\emph{Problems of Information Transmission}, 9(3):265--266, 1973.

\bibitem{Rissanen1978}
J.~Rissanen,
``Modeling by shortest data description,''
\emph{Automatica}, 14(5):465--471, 1978.

\bibitem{Hutter2000}
M.~Hutter,
``A theory of universal artificial intelligence based on algorithmic
complexity,'' arXiv:cs/0004001, 2000.

\bibitem{Schmidhuber2002}
J.~Schmidhuber,
``The Speed Prior: A new simplicity measure yielding near-optimal
computable predictions,''
in \emph{Proc.\ 15th Annual Conf.\ on Computational Learning Theory (COLT)},
pp.~216--228, Springer, 2002.

\bibitem{Hutter2005}
M.~Hutter,
\emph{Universal Artificial Intelligence: Sequential Decisions Based on
Algorithmic Probability},
Springer, Berlin, 2005.

\bibitem{Schmidhuber2007}
J.~Schmidhuber,
``G\"odel machines: Fully self-referential optimal universal self-improvers,''
in \emph{Artificial General Intelligence}, pp.~199--226, Springer, 2007.

\bibitem{Legg2007}
S.~Legg and M.~Hutter,
``Universal intelligence: A definition of machine intelligence,''
\emph{Minds and Machines}, 17(4):391--444, 2007.

\bibitem{Legg2008}
S.~Legg,
\emph{Machine Super Intelligence},
PhD thesis, Universit\`a della Svizzera italiana, 2008.

\bibitem{Grunwald2007}
P.~D.\ Gr\"unwald,
\emph{The Minimum Description Length Principle},
MIT Press, 2007.

\bibitem{Hutter2009Feature}
M.~Hutter,
``Feature reinforcement learning: Part I. Unstructured MDPs,''
\emph{Journal of Artificial General Intelligence}, 1:3--24, 2009.

\bibitem{Veness2011}
J.~Veness, K.~S.~Ng, M.~Hutter, W.~Uther, and D.~Silver,
``A Monte-Carlo AIXI approximation,''
\emph{Journal of Artificial Intelligence Research}, 40:95--142, 2011.

\bibitem{Ghavamzadeh2015}
M.~Ghavamzadeh, S.~Mannor, J.~Pineau, and A.~Tamar,
``Bayesian reinforcement learning: A survey,''
\emph{Foundations and Trends in Machine Learning}, 8(5--6):359--483, 2015.

\bibitem{Orseau2016}
L.~Orseau and S.~Armstrong,
``Safely interruptible agents,''
in \emph{Proc.\ 32nd Conf.\ on Uncertainty in Artificial Intelligence (UAI)},
2016.

\bibitem{Leike2016}
J.~Leike and M.~Hutter,
``On the computability of Solomonoff induction and knowledge-seeking,''
in \emph{Proc.\ 27th International Conf.\ on Algorithmic Learning Theory (ALT)},
pp.~364--378, Springer, 2016.

\bibitem{LiVitanyi2019}
M.~Li and P.~Vit\'anyi,
\emph{An Introduction to Kolmogorov Complexity and Its Applications},
4th ed., Springer, 2019.

\bibitem{Schrittwieser2020}
J.~Schrittwieser, I.~Antonoglou, T.~Hubert, et al.,
``Mastering Atari, Go, Chess and Shogi by planning with a learned model,''
\emph{Nature}, 588:604--609, 2020.

\bibitem{Xie2022}
S.~M.\ Xie, A.~Raghunathan, P.~Liang, and T.~Ma,
``An explanation of in-context learning as implicit Bayesian inference,''
\emph{ICLR}, 2022.

\bibitem{Muller2022}
S.~M\"uller, N.~Hollmann, S.~P.\ Arango, T.~Grabocka, and F.~Hutter,
``Transformers can do Bayesian inference,''
\emph{ICLR}, 2022.

\bibitem{Lotfi2022}
S.~Lotfi, M.~Finzi, S.~Kapoor, A.~Potapczynski, M.~Goldblum, and A.~G.\ Wilson,
``PAC-Bayes compression bounds so tight that they can explain generalization,''
\emph{NeurIPS}, 2022.

\bibitem{Deletang2024}
G.~Del\'etang, A.~Ruoss, P.~A.\ Ortega, et al.,
``Language modeling is compression,''
\emph{ICLR}, 2024.

\end{thebibliography}
\endgroup

\end{document}
